%% CV — Mahima Prajapati
%% Targeted for: Additive Engineering Internship/Co-Op — Freeform (Hawthorne, CA)
%% Compile with: pdflatex main_freeform.tex

\documentclass[11pt,a4paper]{moderncv}
\moderncvstyle{banking}
\moderncvcolor{blue}

\usepackage[utf8]{inputenc}
\usepackage{mathptmx}
\usepackage{hyperref}
\hypersetup{
    colorlinks=true,
    linkcolor=blue,
    filecolor=magenta,
    urlcolor=blue,
    pdftitle={Mahima Prajapati - CV},
    pdfpagemode=FullScreen,
}
\usepackage[scale=0.77]{geometry}
\usepackage{import}

% personal data
\name{Mahima}{Prajapati}
\address{Blacksburg, Virginia, USA}{}{}
\email{mahimap@vt.edu}
\extrainfo{\href{https://linkedin.com/in/mahimaprajapati}{LinkedIn}}

\begin{document}

\makecvtitle

% ============================================================
%     PROFILE STATEMENT
% ============================================================

\vspace{6pt}
\small{PhD candidate in Aerospace Engineering at Virginia Tech specializing in metal additive manufacturing process simulation. Expert in Abaqus-based thermo-mechanical FEA, with research focused on predicting deformations and distortions during laser wire directed energy deposition (DED) and additive friction stir deposition (AFSD) — bridging computational modeling with manufacturing outcomes. Published AIAA SciTech researcher (2026); second paper in preparation. Brings MATLAB scripting, hands-on simulation-driven process insight, and a creative, fast-paced problem-solving approach directly aligned with Freeform's mission to deploy software-defined, autonomous metal 3D printing at industrial scale.}

% ============================================================
%     CORE COMPETENCIES
% ============================================================

\section{Core Competencies}
\vspace{1pt}
\begin{itemize}

\item \textbf{Metal Additive Manufacturing Process Simulation:} Abaqus (expert) — thermo-mechanical FEA of DED and AFSD processes; predicting deformations, residual stress, and distortions in metal AM; bridging computational models to real manufacturing outcomes.
\vspace{1pt}

\item \textbf{Engineering Analysis \& Scripting:} MATLAB — toolpath analysis, numerical methods, data processing, simulation scripting; first-principles approach to evaluating process parameters and part quality.
\vspace{1pt}

\item \textbf{CAD / Design for Manufacturing:} SolidWorks, CATIA — 3D mechanical design, assembly modeling, and design-for-manufacturability (DFM) considerations for additively manufactured components.
\vspace{1pt}

\item \textbf{Structural \& Aeroelastic Analysis:} Abaqus, ANSYS — static and dynamic structural analysis; MS thesis on flutter of aircraft wings with distributed electric propulsors; published at AIAA SciTech 2026.
\vspace{1pt}

\item \textbf{Technical Communication \& Research:} AIAA-published author; experienced instructor to 50+ engineering students; strong written and verbal communication of complex technical concepts; LaTeX documentation.

\end{itemize}

% ============================================================
%     EDUCATION
% ============================================================

\section{Education}
\vspace{1pt}
\begin{itemize}

\item{\cventry{Jan 2025--Present}{Doctor of Philosophy --- Aerospace \& Ocean Engineering}{Virginia Tech}{Blacksburg, VA}{}{\vspace{1pt}
Research: Metal additive manufacturing process simulation; Abaqus FEA of laser wire DED and AFSD; predicting thermo-mechanical deformations and distortions to close the gap between simulation and physical manufacturing.
}}

\vspace{3pt}

\item{\cventry{Aug 2021--Dec 2024}{Master of Science --- Aerospace Engineering}{Virginia Tech}{Blacksburg, VA}{}{\vspace{1pt}
Thesis: ``Aeroelastic Analysis of a Uniform Cantilever Wing with Distributed Electric Propulsors with Rotary Inertia.'' Published and presented at AIAA SciTech Forum 2026. \href{https://doi.org/10.2514/6.2026-1649}{DOI: 10.2514/6.2026-1649}
}}

\vspace{3pt}

\item{\cventry{2016--2020}{Bachelor of Engineering --- Mechanical Engineering}{M.H.\ Saboo Siddik College of Engineering}{Mumbai, India}{}{}}

\end{itemize}

% ============================================================
%     PROFESSIONAL EXPERIENCE
% ============================================================

\section{Professional Experience}
\vspace{3pt}
\begin{itemize}

\item{\cventry{Jan 2025--Present}{PhD Researcher --- Metal Additive Manufacturing}{Virginia Tech, Kevin T. Crofton Dept. of Aerospace \& Ocean Engineering}{Blacksburg, VA}{}{\vspace{1pt}
\begin{itemize}
    \item Developing Abaqus-based thermo-mechanical FEA simulations to predict deformations and distortions during metal AM processes (DED, AFSD); iterating models against experimental data to improve predictive accuracy.
    \vspace{1pt}
    \item Investigating process--structure--property relationships in metal AM; applying simulation outputs to inform process parameter selection and part quality evaluation.
    \vspace{1pt}
    \item Preparing manuscript on thrust effects in aeroelastic analysis of DEP wings (SciTech 2027); journal paper forthcoming.
\end{itemize}}}

\vspace{3pt}

\item{\cventry{Aug 2021--Dec 2024}{Graduate Researcher \& Teaching Assistant --- Aeroelasticity}{Virginia Tech}{Blacksburg, VA}{}{\vspace{1pt}
\begin{itemize}
    \item Developed aeroelastic models for aircraft wings with distributed electric propulsors; conducted flutter speed and frequency analysis; published results at AIAA SciTech 2026.
    \vspace{1pt}
    \item Delivered technical guidance and lab instruction to 50+ undergraduate aerospace engineering students per semester.
    \vspace{1pt}
    \item Competed in Virginia Tech GPSS, earning \textbf{First Place --- Innovation \& Technology Advancement} for poster presentation of aeroelastic flutter research (\$275 award).
\end{itemize}}}


\end{itemize}

% ============================================================
%     PUBLICATIONS
% ============================================================

\section{Publications}
\vspace{1pt}
\begin{itemize}
\item Prajapati, M.~D., and Kapania, R.~K. (2026). ``Aeroelastic Analysis of a Uniform Cantilever Wing with Distributed Electric Propulsors with Rotary Inertia.'' \textit{AIAA SciTech Forum 2026}. \href{https://doi.org/10.2514/6.2026-1649}{DOI: 10.2514/6.2026-1649}
\item Prajapati, M.~D., and Kapania, R.~K. (2027). ``Effect of Thrust on Aeroelastic Analysis of DEP Wings with Rotary Inertia.'' \textit{AIAA SciTech Forum 2027}. (In preparation; journal paper forthcoming.)
\end{itemize}

% ============================================================
%     HONORS AND AWARDS
% ============================================================

\section{Honors and Awards}
\vspace{1pt}
\begin{itemize}
\item{\textbf{First Place --- Innovation \& Technology Advancement} -- Virginia Tech Graduate \& Professional Student Symposium (GPSS); \$275 monetary award.}
\item{\textbf{Tolson \& Nerney Fellowship} -- Kevin T. Crofton Dept.\ of Aerospace \& Ocean Engineering, Virginia Tech (2025); \$11,750.}
\end{itemize}

% ============================================================
%     REFERENCES
% ============================================================

\section{References}
\vspace{1pt}
\begin{itemize}
\item{Available upon request.}
\end{itemize}

\end{document}
